% Analyse
\section{Analyse}\label{sec:analyse}
Es gibt auch die Möglichkeit, zwei Bilder direkt nebeneinander darzustellen:

  \begin{figure} [htbp]
    \centering
    \begin{subfigure}[b]{0.45\textwidth}
        \includegraphics[width=\textwidth]{resources/images/amazon_dev_console.png}
        \caption{Beispielbild4}
        \label{fig:bild4}
    \end{subfigure}
    ~
    \begin{subfigure}[b]{0.45\textwidth}
        \includegraphics[width=\textwidth]{resources/images/amazon_dev_console.png}
        \caption{Beispielbild4}
        \label{fig:bild5}
    \end{subfigure}
    \caption{Zusammenstellung mehrerer Bilder}\label{fig:bild6}
\end{figure}


Um referenzieren auf Kapitel innerhalb deiner Arbeit zu erleichtern, ist in der \lstinline{resources/settings/macros.tex} folgender Befehl hinterlegt: \lstinline!\vgl{String_Kapitelname}{String_label}!. Bei dieser Nutzung \lstinline!\vgl{Einleitung}{sec:einleitung}! wird dies so gerendert: \enquote{\vgl{Einleitung}{sec:einleitung}}. Das erspart dir möglicherweise Zeit beim schreiben.

Zusätzlich kannst du \lstinline!\Vgl! nutzen, um den Vergleich groß zu schreiben, wenn dieser hinter einem Punkt steht. \Vgl{Analyse}{sec:analyse}.
